\chapter{Aims and Context}
A solid-state MASER operable at room temperature and zero magnetic field has
recently been demonstrated in a compact, portable form factor. While this
represents a significant step toward practical implementation, low energy
efficiency, specifically the limited output power relative to the high optical
pump input, continues to constrain real-world adoption of the technology.

This project focuses on the design and optimisation of dielectric resonators
used within the MASER cavity, with the aim of enhancing the Purcell factor and
thereby improving masing efficiency. A genetic algorithm (GA) is employed in
conjunction with COMSOL Multiphysics to evolve resonator geometries that
maximise the Purcell factor by optimising quality factor ($Q_f$) and mode volume
($V_m$). The hypothesis is that irregular resonator shapes can offer superior
confinement and lower radiative loss than traditional cylindrical designs, even
when subject to practical material constraints.

This research is situated within the broader context of low-noise microwave
amplification, a technology with wide-reaching applications in quantum
information systems, deep-space communication, and highly sensitive
spectroscopy. Improving MASER efficiency at room temperature has the potential
to remove key barriers to deployment in industry. As quantum and communication
technologies continue to develop, efficient room-temperature MASERs could form a
foundational component in next-generation sensing and amplification platforms.
